% GENERATED FILE. Edit canonical lessons, then run npm run generate:books.
% canonical-source-hash: fnv1a64:ab047bdebc422c90
% canonical-lessons: RU-C171-pevets, RU-C171-veterinar, RU-C171-sudya, RU-C171-soldat, RU-C171-stroitel

\chapter{Singer, Vet, Judge, Soldier, Builder}
\label{ch:ru-singer-vet-judge-soldier-builder}
\hlchaptermodality{171}

\begin{chapteropening}
\textbf{By the end of this chapter:} \emph{I can say a singer, a vet, a judge, a soldier and a builder.}

Complete the last lesson of chapter 171: I can say a singer, a vet, a judge, a soldier and a builder.
\end{chapteropening}

\section[\texorpdfstring{\ru{певец} (pevéts)}{певец (pevéts)}]{\ru{певец} (pevéts) — a singer}

\label{lesson:RU-C171-pevets}



\begin{quote}
Before the new one: say the Russian for a bill; an account, then the Russian for cash.
\end{quote}

\subsection*{You'll want to know: \ru{певец}}
\textbf{\ru{певец}} (\emph{pevéts}) — \textquotedblleft{}a singer\textquotedblright{}.

\begin{cousinweb}
\textbf{\ru{певец}} comes from \emph{\ru{петь}}, \textquotedblleft{}to sing\textquotedblright{} (Vasmer), with \emph{-\ru{ец}}, \textquotedblleft{}one who does it\textquotedblright{}.
\end{cousinweb}

\begin{grammarlens}[title={What you've built}]
One more word for this chapter.
\end{grammarlens}

\subsection*{Your turn}
\begin{itemize}
\raggedright
  \item \emph{Say it:} \emph{pevéts}
  \item \emph{Say it:} \emph{pevéts}, once more
  \item \emph{Say it:} say \emph{pevéts}
  \item \emph{Recall:} say the Russian for a bill; an account, then the Russian for cash, then say \emph{pevéts} again
\end{itemize}

\begin{tcolorbox}[breakable,colback=teal!4,colframe=teal!35!black,title={Before you move on}]
What does \ru{певец} mean? (\textquotedblleft{}A singer\textquotedblright{}.) Say it once more.
\end{tcolorbox}
\section[\texorpdfstring{\ru{ветеринар} (veterinár)}{ветеринар (veterinár)}]{\ru{ветеринар} (veterinár) — a vet}

\label{lesson:RU-C171-veterinar}



\begin{quote}
Before the new one: say the Russian for cash, then the Russian for a singer.
\end{quote}

\subsection*{You'll want to know: \ru{ветеринар}}
\textbf{\ru{ветеринар}} (\emph{veterinár}) — \textquotedblleft{}a vet\textquotedblright{}.

\begin{cousinweb}
\textbf{\ru{ветеринар}} comes, as German \emph{Veterinär} does, from Latin \emph{veterīnārius}, from \emph{veterīnus}, \textquotedblleft{}a beast of burden\textquotedblright{} (Vasmer).
\end{cousinweb}

\begin{grammarlens}[title={What you've built}]
One more word for this chapter.
\end{grammarlens}

\subsection*{Your turn}
\begin{itemize}
\raggedright
  \item \emph{Say it:} \emph{veterinár}
  \item \emph{Say it:} \emph{veterinár}, once more
  \item \emph{Say it:} say \emph{veterinár}
  \item \emph{Recall:} say the Russian for cash, then the Russian for a singer, then say \emph{veterinár} again
\end{itemize}

\begin{tcolorbox}[breakable,colback=teal!4,colframe=teal!35!black,title={Before you move on}]
What does \ru{ветеринар} mean? (\textquotedblleft{}A vet\textquotedblright{}.) Say it once more.
\end{tcolorbox}
\section[\texorpdfstring{\ru{судья} (sudyá)}{судья (sudyá)}]{\ru{судья} (sudyá) — a judge; a referee}

\label{lesson:RU-C171-sudya}



\begin{quote}
Before the new one: say the Russian for a singer, then the Russian for a vet.
\end{quote}

\subsection*{You'll want to know: \ru{судья}}
\textbf{\ru{судья}} (\emph{sudyá}) — \textquotedblleft{}a judge; a referee\textquotedblright{}.

\begin{cousinweb}
\textbf{\ru{судья}} $\leftarrow$ Proto-Slavic \emph{*sǫdĭji}, built on \emph{\ru{суд}}, \textquotedblleft{}judgement, court\textquotedblright{} — inherited (Vasmer).
\end{cousinweb}

\begin{grammarlens}[title={What you've built}]
One more word for this chapter.
\end{grammarlens}

\subsection*{Your turn}
\begin{itemize}
\raggedright
  \item \emph{Say it:} \emph{sudyá}
  \item \emph{Say it:} \emph{sudyá}, once more
  \item \emph{Say it:} say \emph{sudyá}
  \item \emph{Recall:} say the Russian for a singer, then the Russian for a vet, then say \emph{sudyá} again
\end{itemize}

\begin{tcolorbox}[breakable,colback=teal!4,colframe=teal!35!black,title={Before you move on}]
What does \ru{судья} mean? (\textquotedblleft{}A judge; a referee\textquotedblright{}.) Say it once more.
\end{tcolorbox}
\section[\texorpdfstring{\ru{солдат} (soldát)}{солдат (soldát)}]{\ru{солдат} (soldát) — a soldier}

\label{lesson:RU-C171-soldat}



\begin{quote}
Before the new one: say the Russian for a vet, then the Russian for a judge; a referee.
\end{quote}

\subsection*{You'll want to know: \ru{солдат}}
\textbf{\ru{солдат}} (\emph{soldát}) — \textquotedblleft{}a soldier\textquotedblright{}.

\begin{cousinweb}
\textbf{\ru{солдат}} — known from 1515; it came through German \emph{Soldat} or Dutch \emph{soldaat} from Italian \emph{soldato}, \textquotedblleft{}one hired for pay\textquotedblright{}, from \emph{soldare}, \textquotedblleft{}to hire\textquotedblright{} (Vasmer), and \emph{soldo}, \textquotedblleft{}a coin\textquotedblright{}, Latin \emph{solidus}.
\end{cousinweb}

\begin{grammarlens}[title={What you've built}]
One more word for this chapter.
\end{grammarlens}

\subsection*{Your turn}
\begin{itemize}
\raggedright
  \item \emph{Say it:} \emph{soldát}
  \item \emph{Say it:} \emph{soldát}, once more
  \item \emph{Say it:} say \emph{soldát}
  \item \emph{Recall:} say the Russian for a vet, then the Russian for a judge; a referee, then say \emph{soldát} again
\end{itemize}

\begin{tcolorbox}[breakable,colback=teal!4,colframe=teal!35!black,title={Before you move on}]
What does \ru{солдат} mean? (\textquotedblleft{}A soldier\textquotedblright{}.) Say it once more.
\end{tcolorbox}
\section[\texorpdfstring{\ru{строитель} (stroítel)}{строитель (stroítel)}]{\ru{строитель} (stroítel) — a builder}

\label{lesson:RU-C171-stroitel}



\begin{quote}
Before the new one: say the Russian for a judge; a referee, then the Russian for a soldier.
\end{quote}

\subsection*{You'll want to know: \ru{строитель}}
\textbf{\ru{строитель}} (\emph{stroítel}) — \textquotedblleft{}a builder\textquotedblright{}.

\begin{cousinweb}
\textbf{\ru{строитель}} is \emph{\ru{строить}}, \textquotedblleft{}to build\textquotedblright{}, with \emph{-\ru{тель}}, \textquotedblleft{}one who does it\textquotedblright{}. \emph{\ru{строить}} comes from \emph{\ru{строй}}, \textquotedblleft{}order, structure\textquotedblright{}.
\end{cousinweb}

\begin{grammarlens}[title={What you've built}]
One more word for this chapter.
\end{grammarlens}

\subsection*{Your turn}
\begin{itemize}
\raggedright
  \item \emph{Say it:} \emph{stroítel}
  \item \emph{Say it:} \emph{stroítel}, once more
  \item \emph{Say it:} say \emph{stroítel}
  \item \emph{Recall:} say the Russian for a judge; a referee, then the Russian for a soldier, then say \emph{stroítel} again
\end{itemize}

\begin{tcolorbox}[breakable,colback=teal!4,colframe=teal!35!black,title={Before you move on}]
What does \ru{строитель} mean? (\textquotedblleft{}A builder\textquotedblright{}.) Say it once more.
\end{tcolorbox}
